\section{答题模板}

\subsection{单选题}

\textbf{常用方法}：
\[
\begin{array}{c|c}
\text{方法} & \text{适用场景} \\ \hline
\text{直接法} & \text{公式直接可算} \\
\text{特值法} & \text{含参数的抽象问题} \\
\text{排除法} & \text{不易直接求解} \\
\text{数形结合} & \text{函数零点、方程根}
\end{array}
\]

\paragraph{例题}（2024 新课标Ⅰ卷）已知集合 $A = \{x \mid -5 < x^3 < 5\}$，$B = \{-3,-1,0,2,3\}$，则 $A \cap B =$

A. $\{-1,0\}$ \quad B. $\{2,3\}$ \quad C. $\{-3,-1,0\}$ \quad D. $\{-1,0,2\}$

\textbf{解析}：$x^3 \in (-5,5) \Rightarrow x \in (-\sqrt[3]{5},\sqrt[3]{5})$。$\sqrt[3]{5} \approx 1.71$，故 $A \approx (-1.71, 1.71)$。$B$ 中只有 $-1,0$ 在此区间内。\quad 选 A。

\subsection{多选题}

至少 2 个正确选项，全部选对得 5 分，部分选对得 2 分，有错选得 0 分。\textbf{策略：确定的可选，不确定的不选}。

\subsection{填空题}

直接填写最终结果，无中间过程分。注意定义域、多解情况。

\paragraph{例题}（2024 新课标Ⅰ卷）已知 $\cos^2\alpha + 2\sin2\alpha = \dfrac15$，则 $\tan\alpha = \_\_\_\_\_\_$。

\textbf{解析}：化齐次式：
\[
\cos^2\alpha + 4\sin\alpha\cos\alpha = \frac15,\;
1 + 4\tan\alpha = \frac15(1+\tan^2\alpha)
\]
\[
\tan^2\alpha - 20\tan\alpha - 4 = 0,\;
\tan\alpha = 10 \pm 2\sqrt{26}
\]

\subsection{三角恒等变换与解三角形}

\textbf{模板}：边角互化 $\to$ 三角恒等变形 $\to$ 求值/求角

\textbf{核心公式}：
\[
\frac{a}{\sin A} = \frac{b}{\sin B} = \frac{c}{\sin C} = 2R,\;
a^2 = b^2 + c^2 - 2bc\cos A,\;
S = \frac12 bc\sin A
\]

\paragraph{例题}（2023 新课标Ⅰ卷）在 $\triangle ABC$ 中，$A+B=3C$，$2\sin(A-C)=\sin B$。求 $\sin A$。

\textbf{解析}：$A+B=3C$ 且 $A+B+C=\pi \Rightarrow 4C=\pi \Rightarrow C=\frac{\pi}{4}$。
代入 $2\sin(A-C)=\sin B$，得 $\tan A=3$，$\sin A=\frac{3}{\sqrt{10}}=\frac{3\sqrt{10}}{10}$。

\subsection{数列}

\textbf{模板}：确定类型 $\to$ 求通项 $\to$ 求和

\textbf{方法}：公式法、裂项相消、错位相减、分组求和

\paragraph{例题}（2022 新课标Ⅰ卷）记 $b_n = a_{2n}$，$a_{n+1}=a_n+1$（$n$ 奇）或 $a_n+3$（$n$ 偶），$a_1=1$。

\textbf{解析}：$b_{n+1}=b_n+4$，$b_1=2$，$\therefore b_n=4n-2$。
$S_{20}=$ 奇项和 $+$ 偶项和 $=190+200=300$。

\subsection{立体几何}

\textbf{模板}：建系 $\to$ 坐标 $\to$ 向量法求线面角/二面角/距离

\subsection{概率与统计}

\textbf{模板}：确定模型 $\to$ 分布列 $\to$ 期望/方差

\subsection{解析几何}

\textbf{模板}：设直线 $\to$ 联立 $\to$ 韦达定理 $\to$ 几何条件代数化

\subsection{函数与导数}

\textbf{模板}：求定义域 $\to$ 求导 $\to$ 分析单调性 $\to$ 求最值/证明不等式
